\section{A monotone comparison and localized contours}\label{sec:contours}

\subsection{The auxiliary field and the slab proxy}

For a deterministic slab index \(a\ge0\), initialize
\(\Comp^a(z,a)=1\) at every coarse site. For \(k>a\) use the same
physical bad-bin indicators to define
\begin{equation}\label{eq:comparison}
 \Comp^a(z,k)=
 \begin{cases}
 0,&\Bad(z,k)=1,\\
 \displaystyle\max_{1\le s\le5}\min_{d\in A_s}
                     \Comp^a(z+d,k-1),&\Bad(z,k)=0.
 \end{cases}
\end{equation}
A one denotes a healthy comparison bit. This map is monotone in its
arguments. Also introduce the event, identified with its indicator,
\begin{equation}\label{eq:proxy}
 \Proxy^a(z,k)=
 \{\Bad(z,k)=1\}\ \cup\
 \{\Comp^a(z+d,k-1)=0\text{ for some }d\in\Delta\}.
\end{equation}
Every zero of \(\Comp^a(z,k)\) belongs to \(\Proxy^a(z,k)\).
For the physically empty initial syndrome,
Lemma~\ref{lem:predecessor} gives by induction
\begin{equation}\label{eq:domination}
 \begin{split}
 Q_z\text{ occupied at }kT&\ \Longrightarrow\
                         \Comp^0(z,k)=0,\\
 Q_z\text{ visited in }I_k&\ \Longrightarrow\
                         \Proxy^0(z,k).
 \end{split}
\end{equation}
For the induction, an occupied endpoint of a good bin has an occupied
predecessor in each \(A_s\). Each of those bits is zero, so every
minimum in \eqref{eq:comparison} is zero. For the whole-slab assertion,
the additional implication in Lemma~\ref{lem:predecessor} supplies at
least one occupied predecessor in \(\Delta\). A bad bin belongs to
the proxy without any predecessor requirement.

We use the product metric
\begin{equation}\label{eq:productmetric}
 \rho((z,k),(y,\ell))=\max\{\dist_\infty(z,y),|k-\ell|\}
\end{equation}
on \(\Z_M^4\times\Z\), with periodic spatial distance. On
\(\Z^4\times\Z\), use the ordinary spatial sup distance. A spacetime
path has consecutive distinct sites at \(\rho\)-distance at most
\(R\). A window has dates \(a+1,\ldots,a+D_w\), where on the torus
\begin{equation}\label{eq:short-window}
 2RD_w<M/2.
\end{equation}
On the infinite lattice any finite depth is permitted.

\subsection{The deterministic contour input}

We specify the imported combinatorial input so that its use is
separate from all probabilistic assumptions. Swart--Szab\'o--Toninelli
\cite[Definitions 15--21 and Theorem 22]{SST} work with a countable
typed dependence graph. A vertex carries a constant-zero,
constant-one, or nonconstant monotone map. At a nonconstant vertex
there is a finite nonempty set of outgoing predecessors of each
charge. The associated map is the maximum, over charges, of the
minimum of the predecessor values for that charge. Constants have
no outgoing dependence edges. Their Theorem 22 states that a zero
of the maximal trajectory has a rooted connected Toom contour.

For clarity, the relevant contour structure consists of sources,
sinks and internal vertices carrying one of \(\sigma\) charges.
A source emits one edge of every charge and has no incoming edges;
a sink receives one of each and has no outgoing edges; an internal
vertex has one incoming and one outgoing edge of its charge.
The root is a distinguished source. The zero-edge contour consists
of a single vertex serving as source and sink. An embedding has
distinct sink images, also distinct from all other vertex images,
and distinct images of internal vertices of a given charge, counting
the root with those internal vertices. These are the injectivity
conditions of their Definition 17.

In a present contour the sinks are constant-zero sites and all other
vertices are nonconstant sites. Edges leaving the root or an
internal vertex follow the allowed predecessors of their own charge.
An edge leaving another source may follow a predecessor of any
charge. These are the properties of their Definition 21 and
Eq.~(2.16) that we need. No probability assumption occurs in the
contour implication.

In our application a nonconstant vertex at \((z,k)\) has type-\(s\)
predecessors \((z+d,k-1)\), \(d\in A_s\). All these sets are finite
and nonempty, as required by Definition 20 of \cite{SST}. A bad bin
is constant zero. To use a clean boundary, truncate a root's backward
dependency cone at date \(a\), copy the actual maps inside it, and
put constant-one maps at date \(a\) and everywhere outside, including
the future. This finite acyclic recursion has a unique trajectory
at the root; in particular it agrees there with the maximal
trajectory in Theorem 22. Backward induction on the dates identifies
that value with the root of \eqref{eq:comparison}.

For a torus root, lift this finite cone to \(\Z^4\times\Z\).
A section at a fixed date has coordinate widths at most \(2RD_w\).
Condition~\eqref{eq:short-window} makes projection injective on that
section. Thus distinct lifted sinks at the same date are distinct
physical bins. At distinct dates, bins already have disjoint time
indices. Constant-one cutoff sites cannot be contour vertices, so
every sink is a distinct actual failed test inside the cone. On
the infinite lattice the finite cone needs no lift.

\subsection{Five charges, edge count and dependent sinks}

Use the spatial linear forms of Ahn \cite[Eq.~(7.17)]{Ahn}:
\begin{equation}\label{eq:polar}
 \ell_i(d)=-d_i\quad(1\le i\le4),\qquad
 \ell_5(d)=d_1+d_2+d_3+d_4.
\end{equation}
Their sum is identically zero, and for \(d\in A_s\),
\[
 \ell_s(d)\ge\varepsilon_s,\qquad
       (\varepsilon_1,\ldots,\varepsilon_5)=(0,0,0,0,1).
\]
For any \(d\in\Delta\) the respective losses are at most
\(R,R,R,R,4R\). All contour edges change time by \(-1\).
This checks the step restrictions in \cite[Eq.~(5.7)]{SST}, with
five charges and positive aggregate advance
\(\sum_s\varepsilon_s=1\).

Let a contour have \(m\) sinks. The number of sources is also \(m\),
by counting the net flux of any one charge. By
\cite[Lemma 38, Eq.~(5.12)]{SST}, there are equally many edges of every
charge; denote their common number by \(n\).
This equality also follows directly by summing time differences
along the paths of each charge from sources to sinks. Each edge
has time difference \(-1\), while the endpoint sum is the same for
each charge.

For each charge, \(m-1\) edges leave non-root sources. The other
\(n-m+1\) edges have their prescribed type. Sum
\(\ell_s(\psi(w)-\psi(v))\) over all charged edges \((v,w)\) of charge
\(s\), and over \(s\). Internal contributions cancel; the sum at
each source or sink vanishes by \(\sum_s\ell_s=0\). The total is
therefore zero. The lower bounds on increments give
\[
 0\ge (n-m+1)-8R(m-1).
\]
Thus, with \(\alpha=1+8R\),
\begin{equation}\label{eq:edges-sinks}
 n\le\alpha(m-1),\qquad m\ge1+n/\alpha .
\end{equation}
This is the specialization of \cite[Lemma 40]{SST}; the calculation
checks its polar-function hypotheses directly.

For the contour enumeration we use the version
arXiv:2202.10999v3 of \cite{SST}. With \(\sigma=5\) and
\(\tau=\lceil\sigma/2\rceil=3\), Lemma 39, Eq.~(5.15), states
\(N_n\le n^2 4^{6n}|\Delta|^{5n}\). The sharper odd-charge estimate
in the \emph{proof} of that lemma, Eq.~(5.28), gives
\begin{equation}\label{eq:contour-count}
 N_n\le n^2 4^{6n-1}|\Delta|^{5n}\quad(n\ge1),\qquad N_0=1.
\end{equation}
Here \(N_n\) counts rooted embeddings up to their equivalence of
Definition 19, subject to the step restrictions just verified.
Duplicating the last charge in the proof produces six traversals
per \(n\), but records the duplicated spatial choices once.
The factor \(4^{6n-1}\) is consequently a proof estimate, not the
bound stated in the published Lemma 39. Either estimate is
exponential in \(n\) and would suffice below. Restriction to a
deterministic finite cone can only reduce this count.

Dependence of the bad-bin field is handled without importing an
independent-defect stability theorem. Put
\begin{equation}\label{eq:dependence-divisor}
 D_{\mathrm{dep}}=(4(R+3)+1)^4.
\end{equation}
For any specified set of \(m\) distinct physical bins, join two
when they have the same date and spatial distance at most
\(2(R+3)\). Its maximum degree is at most \(D_{\mathrm{dep}}-1\).
Greedy selection gives at least \(m/D_{\mathrm{dep}}\) bins with disjoint
underlying increments. Lemma~\ref{lem:bad-prob} and independence
therefore imply
\begin{equation}\label{eq:joint-sinks}
 \Prb\{\text{all specified bins are bad}\}\le q^{m/D_{\mathrm{dep}}}.
\end{equation}
Distinct dates give disjoint Poisson intervals even at the same
spatial anchor. For equal dates, the selection ensures disjoint
spatial supports. This is where distinct physical sink indices,
rather than distinct vertices only on a covering lattice, matter.

\subsection{Localized spacetime certificates}

\begin{proposition}[Certificate with adjustable cost]\label{prop:witness}
For every \(A>0\) one can choose \(q,B,\eta_0>0\) in the order of
Lemma~\ref{lem:bad-prob}, independent of \(L,a,D_w\), with the
following property. For each site \(b_0=(z,k)\) in a window satisfying
\eqref{eq:short-window}, there are events
\(\mathcal E^a(b_0,r)\), \(r\in\Z_{\ge1}\), such that
\begin{equation}\label{eq:witness-cover}
 \Proxy^a(z,k)\ \Longrightarrow\
       \bigcup_{r\ge1}\mathcal E^a(b_0,r),\qquad
 \Prb(\mathcal E^a(b_0,r))\le e^{-Ar}.
\end{equation}
Each event reads only basic Poisson increments whose pair
(anchor block, slab index) lies in the closed \(\rho\)-ball
\(\Ball_\rho(b_0,r)\). The assertions hold on \(\Z^4\) at every
finite depth, with the same constants. They also hold conditionally
on the entire physical history up to \(aT\).
\end{proposition}
\begin{proof}
If the current bin is bad, use its single failed test. Otherwise a
zero predecessor \((y,k-1)\), with \(\dist_\infty(y,z)\le R\),
is supplied by \eqref{eq:proxy}. When \(k-1=a\) there is no such
zero. At larger depths apply the finite-cone construction above.
Theorem 22 of \cite{SST} produces a contour whose \(m\) sinks are
distinct bad bins and whose edge count obeys
\eqref{eq:edges-sinks}.

Fix \(K=10(R+3)\). For each \(n\ge0\), use radius
\(r=K(n+1)\). Define \(\mathcal E^a(b_0,r)\) to be the union, over
the at most \((2R+1)^4\) predecessor roots and their candidate
contours with \(n\) edges per charge in the deterministic finite
cones, of the events that the named sinks are bad. At \(n=0\)
include the current failed test as well. At other radii use the
empty event. The sink events are deliberately larger than the
events of actual contour presence; no intermediate good-map
conditions need to be read.

A connected contour has \(5n\) edges. Along an undirected walk from
its root every edge changes space by at most \(R\) and time by one,
even if that walk reverses the edge direction. Thus each sink
lies within \(5Rn+R\) spatial units and \(5n+1\) dates of \(b_0\).
Testing that sink adds \(R+3\) spatial units and no further date.
The radius \(K(n+1)\) contains all those increments. This proves
the asserted locality for every candidate in the union.

Equations~\eqref{eq:contour-count}--\eqref{eq:joint-sinks} give
constants \(C_0,C_1\), depending only on \(R\), such that
\begin{equation}\label{eq:witness-prebound}
 \Prb(\mathcal E^a(b_0,K(n+1)))
 \le C_0C_1^n q^{(1+n/\alpha)/D_{\mathrm{dep}}}\qquad(n\ge0).
\end{equation}
The harmless polynomial \(n^2\) has been included in \(C_1^n\).
Choose \(0<q<1\) small enough that both the constant part and the
coefficient of \(n\) in the logarithm of this bound are at most
\(-AK\). Then \eqref{eq:witness-cover} follows. The physical mark
test realizes this \(q\) by choosing \(B\) and then \(\eta_0\).

The union implication follows by selecting the contour furnished
by the deterministic theorem, or the direct failed test. All cones
and all candidate embeddings are specified without conditioning on
their intervening marks. This also proves the measurability claim.
Every sink lies after \(a\); all increments used are independent
of the history through \(aT\). The same proof therefore gives the
conditional bounds. On \(\Z^4\) no quotient identifies sink indices,
so the argument works at every finite date.
\end{proof}

\paragraph{Witness balls that wrap the torus.}
The radius \(K(n+1)\) is not restricted by the period \(M\).
A spatial ball of radius above \(M/2\) wraps the torus and can cover
all its spatial sites. This is harmless. Its meaning is the actual
closed ball in the periodic product metric, and its support is a
set of physical increments, with no repeated copies. Cone
injectivity was used before this enlargement to charge the sinks
of a single candidate; it is not asserted for the enlarged ball.
The volume bound
\[
 |\Ball_\rho(b_0,r)|\le(2r+1)^5
\]
still holds, because taking a quotient can only decrease the number
of spatial sites relative to that bound. Independence below is used
only for balls disjoint as sets on the actual torus times time.
Wrapping balls that overlap are never declared independent.

\subsection{From certificates to a centered path estimate}

\begin{proposition}[Centered path bound]\label{prop:path}
Choose \(A\) in Proposition~\ref{prop:witness} sufficiently large,
and fix the resulting parameters. There exist \(C,\kappa>0\),
independent of \(L,a,D_w,z,k\), such that for every anchor \(b_0\)
in a permissible window and every real \(u\ge0\),
\begin{equation}\label{eq:path}
 \Prb\left\{\begin{array}{l}
 \text{a finite path in }\Proxy^a\text{ starts at }b_0\\
 \text{and reaches }\rho\text{-distance at least }u
 \end{array}\right\}\le Ce^{-\kappa u}.
\end{equation}
All vertices of the path are required to lie in that window.
On \(\Z^4\) the window can be the entire half-space of dates \(k>a\).
The bound is also valid conditionally on the physical history
through \(aT\).
\end{proposition}
\begin{proof}
Take a finite path realizing the event. Assign to each of its
distinct vertices one occurring witness ball from
Proposition~\ref{prop:witness}. There are finitely many assigned balls.
Select disjoint closed balls greedily in decreasing radius. Each
discarded ball meets a selected ball of at least its radius.
Its center, and indeed the whole discarded ball, lies in the
triple enlargement of that selected ball by the triangle inequality.
Thus selected triple balls cover the path.

If consecutive path vertices have covering centers \(c_i,c_j\)
and selected radii \(r_i,r_j\), then
\[
 \rho(c_i,c_j)\le3r_i+3r_j+R.
\]
The graph of these adjacencies connects the covers of the endpoints.
Choose a simple chain of selected centers between those covers,
possibly a single center. Let its length be \(j\) and
\(s=r_1+\cdots+r_j\). Endpoint displacements contribute at most
\(3r_1+3r_j\). Summing the displayed adjacency bounds gives
\[
 u\le6s+R(j-1)\le7s,
\]
because each allowed radius is at least \(K>R\). This geometric
argument uses the product metric itself and remains valid for
wrapping balls.

It remains to count deterministic candidate chains. For fixed
ordered positive radii summing to the integer \(s\), the initial
center lies within \(3r_1\) of the prescribed anchor, so has at most
\((6r_1+1)^5\) choices. Each subsequent center has at most
\((6r_i+6r_{i+1}+2R+1)^5\) choices. These are upper bounds both
in the lattice and on the torus times time. Since \(r_i\ge K>R\),
the product is at most \(e^{C_2s}\) for an absolute \(C_2\).
For example, the initial factor costs at most \(e^{15r_1}\);
the logarithm of each later factor is at most
\(20(r_i+r_{i+1})\). Thus one may take \(C_2=60\).

There are at most \(2^s\) ordered compositions of \(s\). Count only
chains whose original closed balls are disjoint. Their witness
events read disjoint basic increments, in space and time together,
and therefore have joint probability at most \(e^{-As}\).
Taking \(A>C_2+\log2+1\) and summing gives
\[
 \Prb(\text{path event})
 \le\sum_{s\ge\max\{1,\lceil u/7\rceil\}}
        e^{-(A-C_2-\log2)s}\le Ce^{-\kappa u}.
\]
The union is over deterministic chains, so no independence is
claimed for centers selected after observing the marks. For an
unbounded half-space every witnessing path is finite, and the same
countable union gives the estimate. Conditional independence of
increments after \(aT\) proves its conditional form.
\end{proof}

In particular the bound controls temporal paths of auxiliary zeros
as well as spatial paths of visited blocks. This simultaneous use
of space and time is what permits the window iteration.
